\begin{abstract}
Weight-space model merging has emerged as a computationally efficient paradigm for multi-task model consolidation, yet it suffers from severe parameter interference and representation collapse.
In this work, we approach this challenge from first principles of spectral theory and manifold geometry, proposing \textbf{Grassmannian Subspace Consensus Merging (GSC-Merge)}.
GSC-Merge is a partial weight-space model merging framework that targets the major linear projection layers (comprising over 95\% of Transformer block parameters) while keeping lightweight normalization, biases, and embedding parameters task-specific to prevent statistic mismatch across highly disparate visual domains.
Instead of employing heuristic parameter thresholding or coordinate-wise sign voting, we formulate collective task updates in these layers as a joint multi-task parameter space and construct an optimal low-rank projection operator onto a shared Grassmannian manifold via Singular Value Decomposition (SVD).
Our spectral consensus projection minimizes representation distortion under the Frobenius norm, providing a provable mathematical guarantee on the representational drift of task updates by the Eckart-Young-Mirsky Theorem.
In empirical evaluations using a Vision Transformer backbone across four highly conflicting visual classification tasks, GSC-Merge significantly outperforms coordinate-wise heuristic baselines such as TIES-Merging and Sparse Task Arithmetic. Under both task-conditional swapping and truly task-agnostic settings (where task routing is disabled at inference time), GSC-Merge serves as a robust spectral regularizer, dramatically reducing split-sensitivity variance across validation splits while remaining highly competitive with---and matching the performance of---unconstrained few-shot validation tuning, demonstrating superior multi-task generalization while successfully mitigating the overfitting-optimizer paradox in a low-dimensional consensus subspace.
\end{abstract}
